%% Independent methods file. Not woven into main.tex prose.
%% Every quantity described here resolves to a hash-bound evidence entry.

\subsection{The ladder}

The instrument is a short ordered list of assertions, each of which can be
checked on a given machine without asking the original authors anything. In
increasing order of demand:

\begin{enumerate}
\item A citable published record for the work exists.
\item The reference source can be obtained at a fixed, named revision.
\item The runtime that the reference source requires is available on the machine.
\item The published table can be recomputed on the machine.
\end{enumerate}

Three properties make this more useful than a pass-or-fail verdict. It is
\emph{monotone}: a break at one rung makes every rung above it unreachable, so
the informative quantity is the index of the first break and not a count of
failures. It is \emph{machine-relative}: the same work can reach rung two on one
machine and rung four on another, and the ladder records which machine and when
rather than pretending to a universal property. And it is \emph{refutable}: each
rung names the observation that would settle it, so a reader with a different
machine can disagree with a specific line rather than with a conclusion.

Two things are stipulated not to satisfy a rung, because both are common and
both quietly convert a break into a false pass. A different interpreter that
executes a similar language is not the required runtime; source written for one
numerical environment may run, fail, or silently produce different results under
another, and treating the substitute as equivalent is an assumption, not a
reproduction. And a table produced by a re-implementation is not the published
table, however carefully the re-implementation was written.

Corpus availability is deliberately kept off the ladder and tracked as a separate
axis. It does not behave like the rungs: a corpus can be perfectly public,
openly licensed and permanently archived while still being absent from the
machine, and calling that state either reachable or broken loses the distinction
that matters. We record the public status and the local status as two separate
facts.

\subsection{The probe}

The project being audited already carried a recorded status for each of these
questions. That record is not what this study reports, because a recorded status
ages: a runtime can be installed, a corpus can be fetched, a pinned revision can
be moved. Instead a read-only probe re-checks each rung against the live machine
and writes what it observed with a timestamp, and the manuscript binds the probe
output.

The probe resolves the required interpreter and the common substitute on the
executable search path; confirms that the reference source directory exists,
reads its revision identifier, compares it against the recorded pin, and counts
how many of the recorded entry points are present; and checks a list of candidate
locations for the evaluation corpus, including the location named by the
project's own data-root environment variable. It installs nothing, downloads
nothing, and executes no vendor code. Its output records the first broken rung
directly, so the headline of the audit is a number the probe computed rather than
a reading of prose.

Reporting a probe of this kind is itself a design choice worth stating: an audit
whose evidence is a recollection cannot be re-run, and an audit that cannot be
re-run cannot be contradicted.

\subsection{The comparison that does remain reproducible}

The same project contains a fully local simulation study built on an open-source
room-acoustics package~\cite{scheibler2018pyroomacoustics}, and it supplies the
second half of this paper. Sources are localised from time-of-arrival
measurements at a microphone array in which each channel carries an unknown
constant timing offset. Two estimators are compared. The offset-agnostic
estimator solves an ordinary least-squares problem with all offsets fixed at
zero. The offset-aware estimator treats the offsets as unknown parameters with a
prior and returns a maximum a posteriori estimate over source position and
offsets jointly.

The independent variable is the dispersion of the true per-channel offsets across
\NConditions{} conditions, from a condition in which the offsets are exactly zero
up to a condition in which they are large relative to the measurement noise. Each
condition draws \NTrials{} trials, and both estimators are run on the same trial,
so the comparison is paired within trial. The reported quantity is the median
localisation error. That choice is inherited from the project's own display
rule rather than made here, and it matters: the error distribution has a heavy
upper tail, so the mean is pulled by outliers. Means are recorded in the
underlying artifacts and are not deleted, but they are not promoted to the
display quantity.

Each condition is tested with a two-sided Wilcoxon signed-rank
test~\cite{wilcoxon1945individual} from a standard scientific computing
library~\cite{virtanen2020scipy}, paired per trial. Alongside the test we report
the count of trials on which the offset-aware estimator was the better of the
two, because a count of \NTrials{} paired outcomes distinguishes a small
consistent advantage from a large but mixed one in a way a $p$ value does not.
The simulation was run once at a recorded seed and is not re-run here; this study
reads its output rather than regenerating it.

\subsection{Evidence discipline}

No number in this manuscript is typed by hand. Every artifact it rests on,
including the probe output, is recorded with its SHA-256 digest in a manifest
that accompanies the manuscript source. The figure and table code verifies each
digest against the bytes on disk before reading them and aborts on any mismatch,
and it emits both the figures and the numeric macros that the text prints. Even
the corpus size quoted above is extracted from the bound project record rather
than retyped, so the manuscript cannot disagree with the record it cites.
